\section{NFA dan Epsilon Closure}
\label{sec:ch06-nfa}

\subsection{Definisi}

\begin{definisi}[Nondeterministic Finite Automaton]
NFA adalah lima-tupel
\[
M = (Q, \Sigma, \delta, q_0, F),
\]
dengan \(\delta : Q \times (\Sigma \cup \{\varepsilon\}) \to
\mathcal{P}(Q)\). Untuk setiap state dan simbol, NFA dapat memiliki
nol, satu, atau beberapa state tujuan.
\end{definisi}

Nondeterminisme berarti mesin dapat ``menebak'' langkah yang benar.
String diterima jika terdapat setidaknya satu jalur dari \(q_0\) ke
state penerima. NFA tidak lebih kuat daripada DFA, tetapi sering
lebih mudah ditulis dan dikonstruksi dari ekspresi reguler.

\subsection{Epsilon Closure}

\begin{definisi}[Epsilon closure]
Epsilon closure dari himpunan state \(S\), ditulis
\(\varepsilon\text{-closure}(S)\), adalah himpunan state yang dapat
dicapai dari \(S\) melalui nol atau lebih transisi \(\varepsilon\).
\end{definisi}

Epsilon closure dihitung dengan penelusuran graf. Mulai dari \(S\),
ikuti semua transisi \(\varepsilon\) secara rekursif, dan tambahkan
setiap state yang ditemukan. Kompleksitasnya linear terhadap jumlah
transisi.

\begin{contoh}[Epsilon closure]
Misalkan NFA memiliki transisi
\(q_0 \xrightarrow{\varepsilon} q_1\),
\(q_1 \xrightarrow{\varepsilon} q_2\), dan
\(q_2 \xrightarrow{a} q_3\). Maka
\[
\varepsilon\text{-closure}(\{q_0\}) = \{q_0, q_1, q_2\}.
\]
Transisi \(a\) dari \(q_0\) harus mempertimbangkan seluruh state
dalam closure, sehingga tujuan \(q_3\) dapat dicapai.
\end{contoh}

\subsection{Simulasi NFA}

Simulasi NFA menyimpan himpunan state, bukan satu state:
\begin{enumerate}[label=\arabic*.,leftmargin=*]
  \item mulai dengan \(\varepsilon\text{-closure}(\{q_0\})\);
  \item untuk setiap simbol \(a\), hitung himpunan state tujuan dari
        semua state saat ini, lalu ambil epsilon closure-nya;
  \item jika himpunan state kosong, tolak;
  \item setelah masukan habis, terima jika himpunan state beririsan
        dengan \(F\).
\end{enumerate}

Kompleksitas waktu simulasi NFA adalah \(O(n \cdot |Q|^2)\) pada
implementasi sederhana, karena setiap langkah dapat melibatkan
banyak state. Untuk lexer, NFA biasanya dikonversi menjadi DFA agar
simulasi tetap \(O(n)\).

\begin{figure}[H]
\centering
\begin{tikzpicture}[node distance=20mm]
  \node[state, initial] (q0) {\(0\)};
  \node[state, right=of q0] (q1) {\(1\)};
  \node[state, right=of q1] (q2) {\(2\)};
  \node[accepting, right=of q2] (q3) {\(3\)};
  \draw[arrow] (q0) edge[loop above] node[font=\small] {\(a,b\)} (q0);
  \draw[arrow] (q0) -- node[above, font=\small] {\(a\)} (q1);
  \draw[arrow] (q1) -- node[above, font=\small] {\(b\)} (q2);
  \draw[arrow] (q2) -- node[above, font=\small] {\(b\)} (q3);
\end{tikzpicture}
\caption{NFA untuk \((a \mid b)^{*}abb\)}
\label{fig:ch06-nfa-abb}
\end{figure}

NFA pada \figref{fig:ch06-nfa-abb} menerima string yang berakhir
dengan \(abb\). Pada state 0, mesin dapat tetap di state 0 atau
berpindah ke state 1 ketika membaca \(a\). Pilihan ini adalah inti
nondeterminisme.
